\documentclass[11pt]{article}

% Configuración general
\usepackage[margin=1in]{geometry}
\usepackage[T1]{fontenc}
\usepackage[utf8]{inputenc}
\usepackage[spanish,es-nodecimaldot]{babel}
\usepackage{lmodern}
\usepackage{microtype}

% Matemáticas y estructura
\usepackage{amsmath,amssymb,amsthm,mathtools}
\usepackage{mathrsfs}
\usepackage{enumitem}
\usepackage{xcolor}
\usepackage{booktabs}
\usepackage{array}
\usepackage{tikz}
\usepackage{tikz-cd}
\usetikzlibrary{arrows.meta,babel,backgrounds,calc,decorations.pathreplacing,fit,positioning}
\usepackage{placeins}
\usepackage{hyperref}

\definecolor{azulprofundo}{HTML}{173B57}
\definecolor{azulclaro}{HTML}{EAF2F7}
\hypersetup{
  colorlinks=true,
  linkcolor=azulprofundo,
  urlcolor=azulprofundo,
  pdftitle={Teoría emergente desde la autorrecurrencia}
}

\setlength{\parindent}{0pt}
\setlength{\parskip}{0.65\baselineskip}
\setlist{itemsep=0.25em,topsep=0.35em}
\setcounter{tocdepth}{2}
\numberwithin{equation}{section}
\makeatletter
\renewcommand*\l@section{\@dottedtocline{1}{1.5em}{2.8em}}
\makeatother

\newtheoremstyle{prism}
  {0.7em}{0.7em}{\normalfont}{}%
  {\color{azulprofundo}\bfseries}{.}{0.5em}{}
\theoremstyle{prism}
\newtheorem{axioma}{Axioma}
\newtheorem{definicion}{Definición}
\newtheorem{regla}{Regla}
\newtheorem{observacion}{Observación}
\newtheorem{axiomacap}{Axioma}
\newtheorem{defcap}{Definición}

\newcommand{\F}{\mathcal{F}}
\newcommand{\U}{\mathcal{U}}
\newcommand{\N}{\mathbb{N}}
\newcommand{\raiz}{\rho_{\!*}}
\newcommand{\iter}[2]{#1^{[#2]}}
\newcommand{\peq}{\simeq_{\!r}}
\newcommand{\pentails}{\Rightarrow_{\!r}}
\newcommand{\pneg}{\neg_{\!r}}
\newcommand{\inestable}{\bot_{\!r}}
\newcommand{\Proto}[1]{\boldsymbol{\Xi}_{#1}}
\DeclareMathOperator{\Orb}{Orb}
\DeclareMathOperator{\Con}{Con}
\DeclareMathOperator{\per}{per}
\DeclareMathOperator{\gr}{gr}
\DeclareMathOperator{\Cl}{Cl}

\newcommand{\capitulo}[2]{%
  \clearpage
  \phantomsection
  \addcontentsline{toc}{part}{#1: #2}%
  \begin{center}
    {\color{azulprofundo}\Large\bfseries #1\par}
    \vspace{0.35em}
    {\huge\bfseries #2\par}
  \end{center}
  \vspace{1.2em}
}

\title{\textcolor{azulprofundo}{\bfseries
Teoría emergente desde la autorrecurrencia}\\[0.4em]
\large Una propuesta protoaxiomática}
\author{}
\date{}

\begin{document}

\maketitle
\vspace{-2.5em}

\begin{abstract}
Se propone la autorrecurrencia como mecanismo primitivo para la emergencia de
identidad, estructura y relaciones. A partir de una operación de autoaplicación
se introducen las nociones de protoforma, universo emergente, dimensión,
distancia, aritmética y topología protoestructurales. El objetivo del texto es
organizar estas ideas en un lenguaje matemático uniforme y desarrollar sus
primeras consecuencias fundacionales. Las relaciones de
equivalencia, contención estable y convergencia se consideran primitivas mientras
no se especifique un modelo concreto que las realice.
\end{abstract}

\newpage
\begingroup
\small
\setlength{\parskip}{0pt}
\tableofcontents
\endgroup
\newpage

\phantomsection
\addcontentsline{toc}{part}{Introducción y marco protoaxiomático}

\section{Preliminares y convención de notación}

Sea $\F$ una clase de \emph{formas} dotada de una operación parcial de
aplicación
\[
  (\varphi,\psi)\longmapsto \varphi(\psi).
\]
Cuando las aplicaciones correspondientes estén definidas, la iteración
autorrecursiva de $\varphi$ se denota por
\begin{equation}
  \iter{\varphi}{0}=\varphi,
  \qquad
  \iter{\varphi}{n+1}=\varphi\bigl(\iter{\varphi}{n}\bigr),
  \qquad n\in\N.
  \label{eq:iteracion}
\end{equation}
La relación $\varphi\peq\psi$ significa que ambas formas son
estructuralmente equivalentes bajo recurrencia. No se exige que sean
literalmente el mismo objeto.

\begin{center}
\renewcommand{\arraystretch}{1.2}
\begin{tabular}{>{$}c<{$} p{0.67\textwidth}}
\toprule
\text{Símbolo} & \textbf{Significado} \\
\midrule
\F & Clase de formas sobre la que se define la aplicación. \\
\iter{\varphi}{n} & Iteración autorrecursiva de orden $n$. \\
\varphi\peq\psi & Equivalencia estructural bajo recurrencia. \\
\varphi\pentails\psi & $\psi$ conserva de manera estable la estructura de $\varphi$. \\
\Orb(\varphi) & Órbita autorrecursiva generada por $\varphi$. \\
\bottomrule
\end{tabular}
\end{center}

\section{Protoaxioma fundamental}

\setcounter{axioma}{-1}
\begin{axioma}[Existencia emergente]
Existe una raíz autorrecursiva distinguida $\raiz\in\F$, entendida como el
generador primitivo de la teoría, tal que
\begin{equation}
  \raiz(\raiz)\peq\raiz.
  \label{eq:raiz}
\end{equation}
La raíz se determina mediante su propia acción y no mediante una estructura
externa previamente dada.
\end{axioma}

La autorrecurrencia se adopta como el mecanismo mínimo e irreductible por el
que puede emerger una estructura estable. La unicidad de $\raiz$ se refiere a
su papel fundacional, no a que todas las formas de $\F$ sean idénticas.

\section{Protológica emergente}

La lógica de la teoría no se presupone como una estructura completamente
formada. Se expresa inicialmente mediante las siguientes relaciones:

\begin{description}[style=nextline,leftmargin=1.8em]
  \item[Afirmación autorrecursiva.]
  $\mathsf{A}_r(\varphi)$ afirma que la recurrencia de $\varphi$ preserva su
  identidad; en particular,
  \[
    \mathsf{A}_r(\varphi)
    \quad\Longleftrightarrow\quad
    \varphi(\varphi)\peq\varphi.
  \]

  \item[Negación protoformal.]
  $\pneg\varphi$ indica que la órbita de $\varphi$ no alcanza una forma estable.
  No representa necesariamente la negación clásica.

  \item[Implicación emergente.]
  Se escribe $\varphi\pentails\psi$ cuando $\psi$ contiene o reproduce de
  manera estable la estructura recurrente de $\varphi$.
\end{description}

En este nivel no se impone el principio del tercero excluido. Puede existir un
estado intermedio en el que una recurrencia todavía no se haya estabilizado ni
haya mostrado que no puede hacerlo.

\section{Protoformas y estabilidad}

\begin{definicion}[Protoforma]
Una forma $\varphi\in\F$ es una \emph{protoforma} si su órbita autorrecursiva
conserva una identidad estructural. Es decir, existen $N\in\N$ y
$p\in\N_{>0}$ tales que
\begin{equation}
  \iter{\varphi}{n+p}\peq\iter{\varphi}{n}
  \qquad\text{para todo }n\geq N.
  \label{eq:protoforma}
\end{equation}
\end{definicion}

El caso particular $p=1$ describe un punto fijo estructural:
\[
  \varphi(\varphi)\peq\varphi.
\]
De este modo, la afirmación de que una protoforma es ``su propia causa y
consecuencia'' se interpreta como cierre de su órbita, no como causalidad en
sentido temporal.

La condición de periodicidad eventual puede visualizarse como un diagrama
conmutativo. Si $T_{\varphi}(x)=\varphi(x)$, las flechas horizontales de la
Figura~\ref{fig:diagrama-protoforma} representan aplicaciones sucesivas de
$T_{\varphi}$, mientras que las flechas verticales identifican estados
equivalentes separados por un período $p$.

\begin{figure}[htbp]
  \centering
  \begin{tikzcd}[
    column sep=3.3em,
    row sep=3.2em,
    cells={nodes={font=\small}}
  ]
    \iter{\varphi}{N}
      \arrow[r,"T_{\varphi}"]
      \arrow[d,dashed,"{\simeq_r}"']
    & \iter{\varphi}{N+1}
      \arrow[r,"T_{\varphi}"]
      \arrow[d,dashed,"{\simeq_r}"']
    & \cdots
      \arrow[r,"T_{\varphi}"]
      \arrow[d,dashed,"{\simeq_r}"']
    & \iter{\varphi}{N+p}
      \arrow[d,dashed,"{\simeq_r}"']
    \\
    \iter{\varphi}{N+p}
      \arrow[r,"T_{\varphi}"']
    & \iter{\varphi}{N+p+1}
      \arrow[r,"T_{\varphi}"']
    & \cdots
      \arrow[r,"T_{\varphi}"']
    & \iter{\varphi}{N+2p}
  \end{tikzcd}
  \caption{Diagrama conmutativo de la periodicidad eventual. Cada cuadrado
  expresa que aplicar $T_{\varphi}$ es compatible con la equivalencia
  recurrente $\simeq_r$.}
  \label{fig:diagrama-protoforma}
\end{figure}

Otra manera de entender la misma situación consiste en separar una fase
transitoria y un atractor recurrente. La Figura~\ref{fig:atractor-protoforma}
muestra cómo las primeras iteraciones pueden cambiar de estructura hasta entrar
en un ciclo estable; desde ese momento, la evolución no tiene que permanecer
inmóvil, pero sí repite su patrón tras $p$ aplicaciones.

\begin{figure}[htbp]
  \centering
  \begin{tikzpicture}[
    >=Stealth,
    every node/.style={font=\small},
    transitorio/.style={
      circle,draw=azulprofundo!65,fill=white,thick,minimum size=8.5mm
    },
    estable/.style={
      circle,draw=azulprofundo,fill=azulclaro,very thick,minimum size=10mm
    },
    flujo/.style={->,thick,draw=azulprofundo!75},
    ciclo/.style={->,very thick,draw=azulprofundo}
  ]
    \begin{scope}[on background layer]
      \fill[azulclaro!38] (2.65,0) circle[radius=2.15];
      \draw[azulprofundo!28,dashed,thick]
        (2.65,0) circle[radius=2.15];
      \draw[azulprofundo!16,densely dotted]
        (2.65,0) circle[radius=1.55];
    \end{scope}

    \node[transitorio] (x0) at (-4.8,0.15) {$\iter{\varphi}{0}$};
    \node[transitorio] (x1) at (-3.65,0.85) {$\iter{\varphi}{1}$};
    \node (xd) at (-2.65,0.55) {$\cdots$};
    \node[transitorio] (xn) at (-1.55,1.2) {$\iter{\varphi}{N-1}$};

    \node[estable] (c0) at (1.25,1.35) {$\iter{\varphi}{N}$};
    \node[estable] (c1) at (3.95,1.35) {$\iter{\varphi}{N+1}$};
    \node[estable] (cd) at (4.35,-0.85) {$\cdots$};
    \node[estable] (cp) at (1.05,-0.95) {$\iter{\varphi}{N+p-1}$};
    \node[align=center,text=azulprofundo] at (2.65,0.05)
      {\bfseries atractor recurrente\\[-0.15em]período $p$};

    \draw[flujo] (x0) to[bend left=13] node[above] {$T_{\varphi}$} (x1);
    \draw[flujo] (x1) -- (xd);
    \draw[flujo] (xd) -- (xn);
    \draw[flujo] (xn) to[bend left=10] node[above] {$T_{\varphi}$} (c0);

    \draw[ciclo] (c0) to[bend left=18] node[above] {$T_{\varphi}$} (c1);
    \draw[ciclo] (c1) to[bend left=16] (cd);
    \draw[ciclo] (cd) to[bend left=18] (cp);
    \draw[ciclo] (cp) to[bend left=22]
      node[left] {$T_{\varphi}$} (c0);

    \node[above=0.55em of x1,text=azulprofundo!80]
      {fase transitoria};
    \draw[decorate,decoration={brace,amplitude=4pt},azulprofundo!55]
      (-5.25,-0.55) -- (-1.15,-0.55)
      node[midway,below=0.45em,text=azulprofundo!75]
      {cambios antes de la estabilización};
  \end{tikzpicture}
  \caption{Representación de una órbita eventualmente periódica. Tras una
  etapa transitoria, las iteraciones entran en un ciclo estructural de período
  $p$.}
  \label{fig:atractor-protoforma}
\end{figure}

\FloatBarrier
\section{Universos emergentes}

\begin{axioma}[Universalidad emergente]
Un universo emergente $\U\subseteq\F$ es una colección no vacía de
protoformas que posee relaciones internas estables y está cerrada bajo las
aplicaciones que se encuentren definidas dentro de ella:
\begin{equation}
  \varphi,\psi\in\U
  \quad\text{y}\quad
  \varphi(\psi)\text{ definida}
  \quad\Longrightarrow\quad
  \varphi(\psi)\in\U.
  \label{eq:universo}
\end{equation}
\end{axioma}

Un universo no presupone dimensión espacial, cantidad ni geometría. Estas
propiedades, si aparecen, deben obtenerse de las relaciones recurrentes entre
sus protoformas.

La Figura~\ref{fig:universo-emergente} ofrece una representación esquemática de
esta idea. La frontera de $\U$ no debe interpretarse como un borde espacial:
simboliza que las relaciones y aplicaciones internas no producen formas ajenas
al universo. En particular, si $\varphi,\psi\in\U$, entonces el resultado
$\varphi(\psi)$ permanece en $\U$ siempre que la aplicación esté definida.

\begin{figure}[htbp]
  \centering
  \begin{tikzpicture}[
    >=Stealth,
    every node/.style={font=\small},
    forma/.style={
      circle,draw=azulprofundo,fill=white,very thick,minimum size=10mm
    },
    resultado/.style={
      rounded corners=3pt,draw=azulprofundo,fill=azulclaro,
      very thick,minimum width=16mm,minimum height=9mm
    },
    operador/.style={
      circle,draw=azulprofundo!80,fill=azulprofundo!8,
      thick,minimum size=8mm
    },
    relacion/.style={->,thick,draw=azulprofundo!62},
    aplicacion/.style={->,very thick,draw=azulprofundo}
  ]
    \begin{scope}[on background layer]
      \fill[azulclaro!34] (0,0) ellipse[x radius=6.15cm,y radius=2.75cm];
      \draw[azulprofundo,very thick]
        (0,0) ellipse[x radius=6.15cm,y radius=2.75cm];
      \draw[azulprofundo!25,dashed]
        (0,0) ellipse[x radius=5.75cm,y radius=2.35cm];
    \end{scope}

    \node[anchor=west,text=azulprofundo] at (-5.7,2.3)
      {$\boldsymbol{\mathcal U}$: universo emergente};

    \node[forma] (phi) at (-4.25,0.8) {$\varphi$};
    \node[forma] (psi) at (-2.25,1.65) {$\psi$};
    \node[forma] (chi) at (0.65,1.7) {$\chi$};
    \node[forma] (eta) at (3.55,1.0) {$\eta$};
    \node[forma] (omega) at (4.15,-1.05) {$\omega$};

    \node[operador] (ap) at (-1.85,0.05) {$\mathrm{ap}$};
    \node[resultado] (res) at (0.25,-1.35) {$\varphi(\psi)$};

    \draw[aplicacion] (phi) to[bend right=12]
      node[below left] {forma activa} (ap);
    \draw[aplicacion] (psi) to[bend left=10]
      node[right] {argumento} (ap);
    \draw[aplicacion] (ap) to[bend right=8]
      node[below left] {aplicación} (res);

    \draw[relacion] (psi) to[bend left=12]
      node[above] {$\pentails$} (chi);
    \draw[relacion] (chi) to[bend left=10]
      node[above] {$\pentails$} (eta);
    \draw[relacion] (eta) to[bend left=15] (omega);
    \draw[relacion] (omega) to[bend left=13]
      node[below] {$\peq$} (res);
    \draw[relacion] (res) to[bend left=16]
      node[below] {$\pentails$} (phi);

    \draw[relacion] (chi) to[loop above,min distance=13mm]
      node[above] {$T_{\chi}$} (chi);
    \draw[relacion] (omega) to[loop right,min distance=12mm]
      node[right] {$T_{\omega}$} (omega);

    \node[align=center,text=azulprofundo!85] at (2.0,-0.15)
      {red interna de\\relaciones estables};
    \node[align=center,text=azulprofundo] at (-1.0,-2.25)
      {\bfseries clausura: ninguna aplicación definida sale de $\mathcal U$};
  \end{tikzpicture}
  \caption{Visualización de un universo emergente. Las protoformas forman una
  red de recurrencias internas y la aplicación de $\varphi$ a $\psi$ produce
  otra forma $\varphi(\psi)$ perteneciente al mismo universo.}
  \label{fig:universo-emergente}
\end{figure}

\FloatBarrier
\section{Dimensión como recurrencia}

\begin{definicion}[Dimensión protoestructural]
Para una protoforma $\varphi$, su dimensión protoestructural es el menor
período estable de su órbita:
\begin{equation}
  \mathcal{D}(\varphi)
  =\min\Bigl\{p\in\N_{>0}:\exists N\in\N\quad
  \forall n\geq N,
  \iter{\varphi}{n+p}\peq\iter{\varphi}{n}\Bigr\}.
  \label{eq:dimension}
\end{equation}
Si la órbita no se estabiliza periódicamente, se escribe
$\mathcal{D}(\varphi)=\infty$.
\end{definicion}

Esta dimensión cuenta ciclos internos de recurrencia; no mide longitudes,
áreas ni volúmenes. La formulación evita que la definición sea trivial por el
mero hecho de que $\iter{\varphi}{0}=\varphi$.

La Figura~\ref{fig:dimension-recurrente} presenta tres ``relojes recurrentes''.
La dimensión es el número mínimo de aplicaciones necesarias para regresar a
una clase estructural equivalente una vez alcanzado el régimen estable. Por
ello, un punto fijo tiene dimensión $1$, mientras que una oscilación entre dos o
tres estados tiene dimensión $2$ o $3$, respectivamente.

\begin{figure}[htbp]
  \centering
  \begin{tikzpicture}[
    >=Stealth,
    every node/.style={font=\small},
    estado/.style={
      circle,draw=azulprofundo,fill=white,very thick,minimum size=11mm
    },
    paso/.style={->,very thick,draw=azulprofundo}
  ]
    % Dimensión uno
    \begin{scope}[shift={(-4.6,0)}]
      \fill[azulclaro!38,rounded corners=7pt]
        (-1.8,-2.0) rectangle (1.8,2.0);
      \draw[azulprofundo!35,rounded corners=7pt,thick]
        (-1.8,-2.0) rectangle (1.8,2.0);
      \node[text=azulprofundo,font=\bfseries] at (0,1.55)
        {$\mathcal D(\varphi)=1$};
      \node[estado] (a) at (0,0.25) {$\iter{\varphi}{N}$};
      \draw[paso] (a) to[loop right,min distance=17mm]
        node[right] {$T_{\varphi}$} (a);
      \node[align=center,text=azulprofundo!85] at (0,-1.35)
        {punto fijo\\retorno en un paso};
    \end{scope}

    % Dimensión dos
    \begin{scope}[shift={(0,0)}]
      \fill[azulclaro!38,rounded corners=7pt]
        (-1.8,-2.0) rectangle (1.8,2.0);
      \draw[azulprofundo!35,rounded corners=7pt,thick]
        (-1.8,-2.0) rectangle (1.8,2.0);
      \node[text=azulprofundo,font=\bfseries] at (0,1.55)
        {$\mathcal D(\varphi)=2$};
      \node[estado] (b0) at (-0.85,0.25) {$\iter{\varphi}{N}$};
      \node[estado] (b1) at (0.85,0.25) {$\iter{\varphi}{N+1}$};
      \draw[paso] (b0) to[bend left=24]
        node[above] {$T_{\varphi}$} (b1);
      \draw[paso] (b1) to[bend left=24]
        node[below] {$T_{\varphi}$} (b0);
      \node[align=center,text=azulprofundo!85] at (0,-1.35)
        {oscilación\\retorno en dos pasos};
    \end{scope}

    % Dimensión tres
    \begin{scope}[shift={(4.6,0)}]
      \fill[azulclaro!38,rounded corners=7pt]
        (-1.8,-2.0) rectangle (1.8,2.0);
      \draw[azulprofundo!35,rounded corners=7pt,thick]
        (-1.8,-2.0) rectangle (1.8,2.0);
      \node[text=azulprofundo,font=\bfseries] at (0,1.55)
        {$\mathcal D(\varphi)=3$};
      \node[estado] (c0) at (0,0.95) {$\iter{\varphi}{N}$};
      \node[estado] (c1) at (1.0,-0.15) {$\iter{\varphi}{N+1}$};
      \node[estado] (c2) at (-1.0,-0.15) {$\iter{\varphi}{N+2}$};
      \draw[paso] (c0) to[bend left=13] (c1);
      \draw[paso] (c1) to[bend left=13] (c2);
      \draw[paso] (c2) to[bend left=13]
        node[left] {$T_{\varphi}$} (c0);
      \node[align=center,text=azulprofundo!85] at (0,-1.45)
        {ciclo triple\\retorno en tres pasos};
    \end{scope}

    \draw[decorate,decoration={brace,amplitude=5pt},azulprofundo!55]
      (-6.4,-2.3) -- (6.4,-2.3)
      node[midway,below=0.5em,text=azulprofundo]
      {la dimensión cuenta fases recurrentes, no direcciones espaciales};
  \end{tikzpicture}
  \caption{Ejemplos de dimensión protoestructural. El valor de
  $\mathcal D(\varphi)$ es la longitud del ciclo estable mínimo, medida en
  aplicaciones de $T_{\varphi}$.}
  \label{fig:dimension-recurrente}
\end{figure}

\FloatBarrier
\section{Métrica emergente}

\begin{definicion}[Protodistancia]
La protodistancia entre $\varphi,\psi\in\F$ es el primer nivel en el que sus
órbitas confluyen estructuralmente:
\begin{equation}
  \delta(\varphi,\psi)
  =\inf\Bigl\{n\in\N:
  \iter{\varphi}{n}\peq\iter{\psi}{n}\Bigr\}.
  \label{eq:distancia}
\end{equation}
Si no existe tal nivel, se define $\delta(\varphi,\psi)=\infty$.
\end{definicion}

La función $\delta$ cuantifica estabilidad compartida, no separación espacial.
Para llamarla métrica en sentido estricto sería necesario demostrar, dentro de
un modelo concreto, positividad, simetría y desigualdad triangular.

La Figura~\ref{fig:protodistancia-confluencia} representa la protodistancia como
un primer instante de sincronización. Las dos órbitas pueden recorrer estados
estructuralmente distintos durante varios niveles; el valor
$\delta(\varphi,\psi)=n_{*}$ se alcanza precisamente cuando sus iteraciones son
equivalentes por primera vez.

\begin{figure}[htbp]
  \centering
  \begin{tikzpicture}[
    >=Stealth,
    every node/.style={font=\small},
    estadoA/.style={
      circle,draw=azulprofundo,fill=azulclaro,very thick,minimum size=10mm
    },
    estadoB/.style={
      circle,draw=azulprofundo!65,fill=white,thick,minimum size=10mm
    },
    evolucionA/.style={->,very thick,draw=azulprofundo},
    evolucionB/.style={->,thick,draw=azulprofundo!62}
  ]
    \begin{scope}[on background layer]
      \fill[azulclaro!48,rounded corners=7pt]
        (2.25,-1.75) rectangle (4.45,1.75);
      \draw[azulprofundo!35,dashed,rounded corners=7pt,thick]
        (2.25,-1.75) rectangle (4.45,1.75);
    \end{scope}

    \node[text=azulprofundo,anchor=east] at (-5.65,0.9)
      {\bfseries órbita de $\varphi$};
    \node[text=azulprofundo!70,anchor=east] at (-5.65,-0.9)
      {\bfseries órbita de $\psi$};

    \node[estadoA] (p0) at (-5.1,0.9) {$\iter{\varphi}{0}$};
    \node[estadoA] (p1) at (-2.4,0.9) {$\iter{\varphi}{1}$};
    \node[estadoA] (p2) at (0.3,0.9) {$\iter{\varphi}{2}$};
    \node[estadoA] (pn) at (3.35,0.9) {$\iter{\varphi}{n_{*}}$};

    \node[estadoB] (q0) at (-5.1,-0.9) {$\iter{\psi}{0}$};
    \node[estadoB] (q1) at (-2.4,-0.9) {$\iter{\psi}{1}$};
    \node[estadoB] (q2) at (0.3,-0.9) {$\iter{\psi}{2}$};
    \node[estadoB] (qn) at (3.35,-0.9) {$\iter{\psi}{n_{*}}$};

    \draw[evolucionA] (p0) -- node[above] {$T_{\varphi}$} (p1);
    \draw[evolucionA] (p1) -- node[above] {$T_{\varphi}$} (p2);
    \draw[evolucionA] (p2) -- node[above] {$\cdots$} (pn);

    \draw[evolucionB] (q0) -- node[below] {$T_{\psi}$} (q1);
    \draw[evolucionB] (q1) -- node[below] {$T_{\psi}$} (q2);
    \draw[evolucionB] (q2) -- node[below] {$\cdots$} (qn);

    \draw[azulprofundo!35,densely dotted,thick]
      (p0) -- node[right] {$\not\simeq_r$} (q0);
    \draw[azulprofundo!35,densely dotted,thick]
      (p1) -- node[right] {$\not\simeq_r$} (q1);
    \draw[azulprofundo!35,densely dotted,thick]
      (p2) -- node[right] {$\not\simeq_r$} (q2);
    \draw[<->,azulprofundo,very thick,dashed]
      (pn) -- node[right,fill=azulclaro!48,inner sep=2pt]
      {$\simeq_r$} (qn);

    \node[align=center,text=azulprofundo,anchor=south] at (3.35,1.78)
      {\bfseries primera confluencia};
    \node[align=center,text=azulprofundo] at (5.35,0)
      {$\boxed{\delta(\varphi,\psi)=n_{*}}$};

    \draw[->,azulprofundo!70,thick] (-5.55,-2.15) -- (4.65,-2.15)
      node[right] {nivel de recurrencia};
    \foreach \x/\etiqueta in {-5.1/0,-2.4/1,0.3/2,3.35/n_{*}} {
      \draw[azulprofundo!70] (\x,-2.05) -- (\x,-2.25)
        node[below] {$\etiqueta$};
    }
  \end{tikzpicture}
  \caption{La protodistancia como primer nivel de confluencia estructural. Las
  órbitas permanecen distinguibles para $n<n_{*}$ y se vuelven equivalentes en
  el nivel $n_{*}$.}
  \label{fig:protodistancia-confluencia}
\end{figure}

\FloatBarrier
\section{Protoaritmética}

Los números protoestructurales representan niveles o tipos de estabilización,
en lugar de cantidades ordinarias. Para $n\in\N_{>0}$, se introduce la clase
\begin{axioma}[Números protoestructurales]
\begin{equation}
  \mathbf{P}_n
  =\bigl\{\varphi\in\F:\mathcal{D}(\varphi)=n\bigr\}.
  \label{eq:protonumero}
\end{equation}
Así, el símbolo $n$ puede emplearse como etiqueta del comportamiento recurrente
de período $n$.
\end{axioma}

Cuando $a$ y $b$ representan acciones recurrentes compatibles, se proponen las
operaciones formales
\begin{align}
  a\oplus b &\coloneqq a(b),
  \\
  a\otimes b &\coloneqq a^{\circ b}
  =\underbrace{a\circ a\circ\cdots\circ a}_{b\text{ composiciones}}.
\end{align}
Estas operaciones no deben identificarse con la suma y el producto usuales sin
probar previamente propiedades como clausura, asociatividad y existencia de
elementos neutros.

La Figura~\ref{fig:protoaritmetica-circuitos} separa los dos mecanismos
operacionales. La operación $\oplus$ introduce una forma como argumento de otra,
mientras que $\otimes$ representa la reiteración de una misma acción. El segundo
panel utiliza una forma auxiliar $x$ para hacer visible el efecto de la
composición repetida.

\begin{figure}[htbp]
  \centering
  \begin{tikzpicture}[
    >=Stealth,
    every node/.style={font=\small},
    entrada/.style={circle,draw=azulprofundo!70,fill=white,thick,minimum size=9mm},
    operador/.style={rounded corners=4pt,draw=azulprofundo,fill=azulclaro,
      very thick,minimum width=11mm,minimum height=10mm},
    salida/.style={rounded corners=4pt,draw=azulprofundo,fill=white,
      very thick,minimum width=17mm,minimum height=10mm},
    flujo/.style={->,very thick,draw=azulprofundo,shorten >=2pt,shorten <=2pt}
  ]
    % Panel de aplicación
    \begin{scope}[shift={(-3.8,0)}]
      \fill[azulclaro!28,rounded corners=8pt]
        (-2.5,-2.0) rectangle (2.5,2.0);
      \draw[azulprofundo!35,rounded corners=8pt,thick]
        (-2.5,-2.0) rectangle (2.5,2.0);
      \node[text=azulprofundo,font=\bfseries] at (0,1.55)
        {aplicación: $a\oplus b$};
      \node[entrada] (b) at (-1.65,0.15) {$b$};
      \node[operador] (opa) at (-0.05,0.15) {$a(\,\cdot\,)$};
      \node[salida] (ab) at (1.65,0.15) {$a(b)$};
      \draw[flujo] (b) -- node[above] {argumento} (opa);
      \draw[flujo] (opa) -- node[above] {acción} (ab);
      \node[align=center,text=azulprofundo!80] at (0,-1.25)
        {dos formas compatibles\\producen una nueva forma};
    \end{scope}

    % Panel de composición reiterada
    \begin{scope}[shift={(3.45,0)}]
      \fill[azulclaro!28,rounded corners=8pt]
        (-3.7,-2.0) rectangle (3.7,2.0);
      \draw[azulprofundo!35,rounded corners=8pt,thick]
        (-3.7,-2.0) rectangle (3.7,2.0);
      \node[text=azulprofundo,font=\bfseries] at (0,1.55)
        {iteración: $a\otimes b=a^{\circ b}$};
      \node[entrada] (x) at (-3.05,0.15) {$x$};
      \node[operador] (a1) at (-1.75,0.15) {$a$};
      \node[operador] (a2) at (-0.35,0.15) {$a$};
      \node (dots) at (0.7,0.15) {$\cdots$};
      \node[operador] (ak) at (1.55,0.15) {$a$};
      \node[salida] (out) at (3.0,0.15) {$a^{\circ b}(x)$};
      \draw[flujo] (x) -- (a1);
      \draw[flujo] (a1) -- (a2);
      \draw[flujo] (a2) -- (dots);
      \draw[flujo] (dots) -- (ak);
      \draw[flujo] (ak) -- (out);
      \draw[decorate,decoration={brace,amplitude=4pt},azulprofundo!60]
        (-2.25,-0.65) -- (2.25,-0.65)
        node[midway,below=0.45em,text=azulprofundo] {$b$ aplicaciones};
      \node[align=center,text=azulprofundo!80] at (0,-1.55)
        {la recurrencia convierte repetición\\en estructura operacional};
    \end{scope}
  \end{tikzpicture}
  \caption{Interpretación operacional de la protoaritmética. A la izquierda,
  $\oplus$ actúa por aplicación; a la derecha, $\otimes$ codifica composición
  recurrente.}
  \label{fig:protoaritmetica-circuitos}
\end{figure}

\FloatBarrier
\section{Prototopología}

\begin{axioma}[Protovecindad]
Para $\varepsilon\in\N_{>0}$, la vecindad recurrente de radio $\varepsilon$
centrada en $\varphi$ es
\begin{equation}
  B_{\varepsilon}(\varphi)
  =\bigl\{\psi\in\F:\delta(\varphi,\psi)<\varepsilon\bigr\}.
  \label{eq:vecindad}
\end{equation}
Dos protoformas son cercanas al nivel $\varepsilon$ si
$\psi\in B_{\varepsilon}(\varphi)$.
\end{axioma}

La cercanía expresa confluencia autorrecursiva. Una topología propiamente dicha
se obtiene cuando estas vecindades generan una familia que contiene a $\F$ y al
vacío, y es estable bajo uniones arbitrarias e intersecciones finitas.

La Figura~\ref{fig:prototopologia-vecindades} muestra una familia anidada de
protovecindades. Las circunferencias son una ayuda visual y no distancias
espaciales: cada frontera agrupa las formas cuya confluencia con $\varphi$
ocurre antes del umbral recurrente indicado.

\begin{figure}[htbp]
  \centering
  \begin{tikzpicture}[
    >=Stealth,
    every node/.style={font=\small},
    forma/.style={circle,draw=azulprofundo,fill=white,thick,minimum size=8.5mm},
    centro/.style={circle,draw=azulprofundo,fill=azulprofundo,
      text=white,very thick,minimum size=10mm}
  ]
    \begin{scope}[shift={(-1.7,0)}]
      \fill[azulclaro!14] (0,0) circle[radius=3.25];
      \fill[azulclaro!24] (0,0) circle[radius=2.35];
      \fill[azulclaro!42] (0,0) circle[radius=1.35];
      \draw[azulprofundo!45,dashed,thick] (0,0) circle[radius=3.25];
      \draw[azulprofundo!62,dashed,thick] (0,0) circle[radius=2.35];
      \draw[azulprofundo,very thick] (0,0) circle[radius=1.35];

      \node[centro] (phi) at (0,0) {$\varphi$};
      \node[forma] (p1) at (0.75,0.45) {$\psi_1$};
      \node[forma] (p2) at (-1.45,1.05) {$\psi_2$};
      \node[forma] (p3) at (1.55,-1.25) {$\psi_3$};
      \node[forma] (p4) at (-2.55,-1.25) {$\psi_4$};
      \node[forma] (p5) at (2.65,-1.75) {$\psi_5$};
      \node[forma,draw=azulprofundo!40,text=azulprofundo!55]
        (out) at (-3.75,2.4) {$\xi$};

      \draw[->,azulprofundo!55,thick] (p1) -- node[above] {$\delta<1$} (phi);
      \draw[->,azulprofundo!45,thick] (p3) to[bend right=13] (phi);
      \draw[->,azulprofundo!35,thick] (p5) to[bend left=15] (phi);

      \node[text=azulprofundo,anchor=west] at (0.95,1.05)
        {$B_{1}(\varphi)$};
      \node[text=azulprofundo!78,anchor=west] at (1.55,1.65)
        {$B_{2}(\varphi)$};
      \node[text=azulprofundo!62,anchor=west] at (2.45,2.45)
        {$B_{3}(\varphi)$};
      \node[text=azulprofundo!55,anchor=east] at (-3.35,3.0)
        {$\xi\notin B_{3}(\varphi)$};
    \end{scope}

    \node[align=left,text=azulprofundo] at (4.15,1.6)
      {$B_{1}(\varphi)\subseteq B_{2}(\varphi)$\\[0.35em]
       $\hphantom{B_{1}(\varphi)}\subseteq B_{3}(\varphi)$};
    \node[align=left,text=azulprofundo!80] at (4.15,0.15)
      {menor $\delta$\\$\Longrightarrow$\\mayor cercanía recurrente};
    \draw[->,very thick,azulprofundo!65] (3.35,-1.15) -- (5.05,-1.15)
      node[midway,below=0.45em] {umbral $\varepsilon$ creciente};
  \end{tikzpicture}
  \caption{Protovecindades anidadas centradas en $\varphi$. La pertenencia a
  $B_{\varepsilon}(\varphi)$ depende del nivel de confluencia recurrente y no
  de una distancia espacial previa.}
  \label{fig:prototopologia-vecindades}
\end{figure}

\FloatBarrier
\section{Metaestructura cíclica}

\begin{definicion}[Metaestructura]
La metaestructura cíclica asociada a $\varphi$ es su órbita:
\begin{equation}
  \mathscr{M}(\varphi)
  \coloneqq\Orb(\varphi)
  =\bigl\{\iter{\varphi}{0},\iter{\varphi}{1},
  \iter{\varphi}{2},\ldots\bigr\}.
  \label{eq:metaestructura}
\end{equation}
\end{definicion}

La secuencia representa el pliegue y despliegue de la autorrecurrencia a través
de múltiples niveles. Si dos términos distintos de la secuencia son equivalentes,
se forma un ciclo; dicho ciclo se interpreta como un nodo emergente de
significado.

La Figura~\ref{fig:metaestructura-pliegue} visualiza este proceso como una cinta
de iteraciones. La sucesión se despliega hacia estados nuevos hasta que una
iteración posterior reproduce la clase estructural de un estado anterior. La
flecha de retorno pliega entonces la cinta y delimita un ciclo con significado
propio.

\begin{figure}[htbp]
  \centering
  \begin{tikzpicture}[
    >=Stealth,
    every node/.style={font=\small},
    estado/.style={circle,draw=azulprofundo,fill=white,very thick,minimum size=10mm},
    cicloestado/.style={circle,draw=azulprofundo,fill=azulclaro,
      very thick,minimum size=10mm},
    avance/.style={->,very thick,draw=azulprofundo}
  ]
    \node[estado] (s0) at (-5.4,0) {$\iter{\varphi}{0}$};
    \node[cicloestado] (s1) at (-3.1,0.9) {$\iter{\varphi}{1}$};
    \node[cicloestado] (s2) at (-0.7,0) {$\iter{\varphi}{2}$};
    \node[cicloestado] (s3) at (1.7,-0.9) {$\iter{\varphi}{3}$};
    \node[cicloestado] (s4) at (4.1,0) {$\iter{\varphi}{4}$};

    \begin{scope}[on background layer]
      \node[fit=(s1)(s2)(s3)(s4),rounded corners=14pt,
        fill=azulclaro!32,draw=azulprofundo!35,dashed,thick,
        inner sep=13mm] (zona) {};
      \fill[azulprofundo!5,rounded corners=12pt]
        (-6.15,-0.55) -- (-3.25,1.55) -- (-0.55,0.65) --
        (1.8,-1.55) -- (4.75,-0.45) -- (4.55,0.45) --
        (1.65,-0.25) -- (-0.85,0.65) -- (-3.0,1.25) -- cycle;
    \end{scope}

    \draw[avance] (s0) to[bend left=8] node[above left] {$T_{\varphi}$} (s1);
    \draw[avance] (s1) to[bend left=8] node[above] {$T_{\varphi}$} (s2);
    \draw[avance] (s2) to[bend left=8] node[below left] {$T_{\varphi}$} (s3);
    \draw[avance] (s3) to[bend left=8] node[below right] {$T_{\varphi}$} (s4);
    \draw[->,azulprofundo,very thick,dashed]
      (s4) to[bend right=55] node[above] {$\iter{\varphi}{4}\peq\iter{\varphi}{1}$} (s1);

    \node[text=azulprofundo!75] at (-5.35,-1.0) {origen};
    \node[text=azulprofundo,font=\bfseries] at (0.55,2.05)
      {pliegue recurrente};

    \node[circle,draw=azulprofundo,double,fill=white,very thick,
      minimum size=14mm] (nodo) at (5.75,-1.1) {$\mathfrak n_{\varphi}$};
    \draw[->,very thick,azulprofundo!70] (s4) to[bend left=18]
      node[right] {síntesis} (nodo);
    \node[align=center,text=azulprofundo,anchor=north] at (5.75,-1.9)
      {nodo emergente\\de significado};

    \draw[decorate,decoration={brace,amplitude=5pt},azulprofundo!55]
      (-5.9,-2.45) -- (4.65,-2.45)
      node[midway,below=0.5em,text=azulprofundo]
      {despliegue de la órbita a través de niveles recurrentes};
  \end{tikzpicture}
  \caption{Metaestructura cíclica como despliegue y pliegue de una órbita.
  La repetición estructural cierra un ciclo y permite tratarlo como un nodo
  emergente $\mathfrak n_{\varphi}$.}
  \label{fig:metaestructura-pliegue}
\end{figure}

\FloatBarrier
\section{Autorrecurrencia, autosimilitud y equidad}

Dentro de la jerarquía conceptual propuesta, las tres nociones cumplen
funciones diferentes:
\begin{enumerate}
  \item la \textbf{autorrecurrencia} proporciona el mecanismo generador;
  \item la \textbf{autosimilitud} describe la persistencia de un patrón entre
  niveles de recurrencia;
  \item la \textbf{equidad} expresa un estado emergente de simetría o balance
  entre formas ya constituidas.
\end{enumerate}
En consecuencia, el orden de dependencia sugerido es
\[
  \boxed{\text{autorrecurrencia}}
  \longrightarrow
  \boxed{\text{autosimilitud}}
  \longrightarrow
  \boxed{\text{equidad}}.
\]
La flecha indica dependencia conceptual, no implicación lógica automática.

\section{Colecciones emergentes}

\begin{axioma}[Protoconjuntismo]
Una protoforma no se considera aislada: se caracteriza también por las formas
con las que mantiene una recurrencia estable. La colección emergente generada
por $\varphi$ es
\begin{equation}
  \Con(\varphi)
  =\bigl\{\psi\in\F:
  \varphi\pentails\psi\quad\text{o}\quad\psi\pentails\varphi\bigr\}.
  \label{eq:conjunto}
\end{equation}
\end{axioma}

En la interpretación original no existe un elemento absolutamente aislado. El
vacío emergente puede entenderse como ausencia de estabilización dentro de una
órbita. Esto no elimina el conjunto vacío $\varnothing$ de la metateoría, pues
este sigue siendo necesario para formular la teoría de conjuntos y la topología
que sirven como lenguaje externo.

La Figura~\ref{fig:coleccion-emergente} representa $\Con(\varphi)$ como una
constelación relacional. Su contorno es únicamente esquemático: una forma
pertenece a la colección no por ocupar una región, sino porque mantiene con
$\varphi$ al menos una relación recurrente estable en alguna dirección.

\begin{figure}[htbp]
  \centering
  \begin{tikzpicture}[
    >=Stealth,
    every node/.style={font=\small},
    centro/.style={
      circle,draw=azulprofundo,fill=azulprofundo,text=white,
      very thick,minimum size=13mm
    },
    miembro/.style={
      circle,draw=azulprofundo,fill=white,very thick,minimum size=10mm
    },
    exterior/.style={
      circle,draw=azulprofundo!38,fill=white,text=azulprofundo!55,
      dashed,thick,minimum size=10mm
    },
    relacion/.style={->,very thick,draw=azulprofundo},
    relacionsec/.style={->,thick,draw=azulprofundo!62}
  ]
    \begin{scope}[on background layer]
      \fill[azulclaro!34]
        plot[smooth cycle,tension=0.75] coordinates {
          (-4.8,0.1) (-3.5,2.35) (-0.7,2.7) (2.25,2.25)
          (3.45,0.3) (2.55,-2.15) (-0.15,-2.55) (-3.4,-2.05)
        };
      \draw[azulprofundo,very thick]
        plot[smooth cycle,tension=0.75] coordinates {
          (-4.8,0.1) (-3.5,2.35) (-0.7,2.7) (2.25,2.25)
          (3.45,0.3) (2.55,-2.15) (-0.15,-2.55) (-3.4,-2.05)
        };
      \draw[azulprofundo!25,dashed,thick]
        plot[smooth cycle,tension=0.75] coordinates {
          (-4.45,0.05) (-3.25,2.0) (-0.7,2.35) (1.95,1.95)
          (3.05,0.25) (2.25,-1.8) (-0.2,-2.2) (-3.1,-1.75)
        };
    \end{scope}

    \node[anchor=west,text=azulprofundo,font=\bfseries]
      at (-4.15,2.1) {$\Con(\varphi)$};

    \node[centro] (phi) at (-0.65,0) {$\varphi$};
    \node[miembro] (p1) at (-3.25,0.9) {$\psi_1$};
    \node[miembro] (p2) at (-2.65,-1.25) {$\psi_2$};
    \node[miembro] (p3) at (0.3,1.55) {$\psi_3$};
    \node[miembro] (p4) at (1.8,-0.9) {$\psi_4$};
    \node[miembro] (p5) at (2.1,0.95) {$\psi_5$};

    \draw[relacion] (phi) to[bend left=12]
      node[above] {$\pentails$} (p1);
    \draw[relacion] (p2) to[bend left=10]
      node[below right] {$\pentails$} (phi);
    \draw[relacion] (phi) to[bend left=14]
      node[right] {$\pentails$} (p3);
    \draw[relacionsec] (p3) to[bend left=14] (phi);
    \draw[relacion] (p4) to[bend left=12]
      node[below] {$\pentails$} (phi);
    \draw[relacion] (phi) to[bend left=10]
      node[above] {$\pentails$} (p5);

    \draw[relacionsec] (p1) to[bend left=13] (p3);
    \draw[relacionsec] (p3) to[bend left=13] (p5);
    \draw[relacionsec] (p2) to[bend right=14] (p4);

    \node[exterior] (xi) at (5.0,1.15) {$\xi$};
    \draw[azulprofundo!38,dashed,thick]
      (3.3,0.85) -- node[above,sloped,text=azulprofundo!58]
      {sin recurrencia estable} (xi);
    \node[align=center,text=azulprofundo!65] at (5.0,0.25)
      {$\xi\notin\Con(\varphi)$};

    \node[rounded corners=5pt,draw=azulprofundo!55,fill=white,
      align=center,inner sep=7pt,text=azulprofundo] at (4.85,-1.35)
      {$\psi\in\Con(\varphi)$\\[0.2em]
       $\Longleftrightarrow$\\[-0.05em]
       $\varphi\pentails\psi$ o $\psi\pentails\varphi$};
  \end{tikzpicture}
  \caption{Colección emergente generada por $\varphi$. Las protoformas
  conectadas con el centro mediante una recurrencia estable pertenecen a
  $\Con(\varphi)$, independientemente del sentido de la relación; una forma
  sin tal vínculo queda fuera de la colección.}
  \label{fig:coleccion-emergente}
\end{figure}

\FloatBarrier

\section{Síntesis axiomática}

La propuesta puede resumirse en los siguientes principios:
\begin{enumerate}[label=\textbf{P\arabic*.}]
  \item \textbf{Existencia emergente:} una entidad interna se reconoce mediante
  un comportamiento autorrecurrente.
  \item \textbf{Protoforma:} la identidad estructural consiste en la estabilidad
  eventual de una órbita.
  \item \textbf{Dimensión:} la dimensión protoestructural es el período mínimo de
  recurrencia estable.
  \item \textbf{Medición:} se miden niveles de recurrencia y confluencia, no
  cantidades espaciales primitivas.
  \item \textbf{Topología:} la cercanía procede de la confluencia entre órbitas.
  \item \textbf{Protológica:} el discernimiento se basa en la preservación y la
  contención estable de estructura.
  \item \textbf{Colecciones:} las colecciones emergen de relaciones recurrentes
  entre formas.
  \item \textbf{Similitud:} la autosimilitud es resonancia estructural entre
  distintos niveles.
  \item \textbf{Equidad:} la equidad es una metaforma de balance recurrente.
\end{enumerate}

\section{Sistema protodeductivo}

Se adoptan tres expresiones básicas:
\[
  \mathsf{A}_r(\varphi),
  \qquad
  \varphi\vdash_r\psi,
  \qquad
  \varphi\peq\psi.
\]
La primera afirma la estabilidad autorrecursiva de $\varphi$; la segunda indica
que la estructura de $\psi$ se obtiene recurrentemente a partir de $\varphi$; y
la tercera expresa convergencia o equivalencia recurrente.

\begin{regla}[Estabilidad]
Si $\varphi(\varphi)\peq\varphi$, entonces
$\mathsf{A}_r(\varphi)$.
\end{regla}

\begin{regla}[Transitividad]
Si $\varphi\vdash_r\psi$ y $\psi\vdash_r\chi$, entonces
$\varphi\vdash_r\chi$.
\end{regla}

\begin{regla}[Equivalencia por recurrencia mutua]
Si $\varphi\vdash_r\psi$ y $\psi\vdash_r\varphi$, entonces
$\varphi\peq\psi$.
\end{regla}

La Figura~\ref{fig:sistema-protodeductivo} organiza las tres reglas como canales
de transformación inferencial. En cada canal, unas premisas recurrentes ingresan
al núcleo de la regla y producen una conclusión cuya información estructural ya
estaba contenida en ellas.

\begin{figure}[htbp]
  \centering
  \begin{tikzpicture}[
    >=Stealth,
    every node/.style={font=\small},
    premisa/.style={
      rounded corners=4pt,draw=azulprofundo!65,fill=white,thick,
      minimum width=37mm,minimum height=11mm,align=center,inner sep=5pt
    },
    nucleo/.style={
      rounded corners=6pt,draw=azulprofundo,fill=azulclaro,very thick,
      minimum width=27mm,minimum height=13mm,align=center
    },
    conclusion/.style={
      rounded corners=4pt,draw=azulprofundo,fill=azulprofundo,
      text=white,very thick,minimum width=31mm,minimum height=11mm,
      align=center,inner sep=5pt
    },
    flujo/.style={->,very thick,draw=azulprofundo,shorten >=2pt,shorten <=2pt}
  ]
    \begin{scope}[on background layer]
      \foreach \y in {2.15,0,-2.15} {
        \fill[azulclaro!23,rounded corners=8pt]
          (-6.1,\y-0.82) rectangle (6.1,\y+0.82);
        \draw[azulprofundo!24,rounded corners=8pt,thick]
          (-6.1,\y-0.82) rectangle (6.1,\y+0.82);
      }
    \end{scope}

    \node[text=azulprofundo,font=\bfseries,anchor=west] at (-5.85,3.35)
      {premisas recurrentes};
    \node[text=azulprofundo,font=\bfseries] at (0,3.35)
      {regla};
    \node[text=azulprofundo,font=\bfseries,anchor=east] at (5.85,3.35)
      {conclusión emergente};
    \draw[->,azulprofundo!55,thick] (-5.8,3.0) -- (5.8,3.0)
      node[midway,fill=white,inner sep=2pt,text=azulprofundo!75]
      {flujo protodeductivo};

    % Canal de estabilidad
    \node[text=azulprofundo,font=\bfseries] at (-5.65,2.15) {R1};
    \node[premisa] (e-prem) at (-3.45,2.15)
      {$\varphi(\varphi)\peq\varphi$};
    \node[nucleo] (e-regla) at (0,2.15)
      {\bfseries estabilidad};
    \node[conclusion] (e-conc) at (4.0,2.15)
      {$\mathsf A_r(\varphi)$};
    \draw[flujo] (e-prem) -- (e-regla);
    \draw[flujo] (e-regla) -- (e-conc);

    % Canal de transitividad
    \node[text=azulprofundo,font=\bfseries] at (-5.65,0) {R2};
    \node[premisa] (t-prem) at (-3.45,0)
      {$\varphi\vdash_r\psi$\\[-0.05em]
       $\psi\vdash_r\chi$};
    \node[nucleo] (t-regla) at (0,0)
      {\bfseries transitividad};
    \node[conclusion] (t-conc) at (4.0,0)
      {$\varphi\vdash_r\chi$};
    \draw[flujo] (t-prem) -- (t-regla);
    \draw[flujo] (t-regla) -- (t-conc);

    % Canal de equivalencia mutua
    \node[text=azulprofundo,font=\bfseries] at (-5.65,-2.15) {R3};
    \node[premisa] (m-prem) at (-3.45,-2.15)
      {$\varphi\vdash_r\psi$\\[-0.05em]
       $\psi\vdash_r\varphi$};
    \node[nucleo] (m-regla) at (0,-2.15)
      {\bfseries recurrencia\\[-0.1em]mutua};
    \node[conclusion] (m-conc) at (4.0,-2.15)
      {$\varphi\peq\psi$};
    \draw[flujo] (m-prem) -- (m-regla);
    \draw[flujo] (m-regla) -- (m-conc);

    \draw[decorate,decoration={brace,amplitude=5pt},azulprofundo!55]
      (-5.95,-3.25) -- (5.95,-3.25)
      node[midway,below=0.5em,text=azulprofundo]
      {preservación, composición y reconocimiento de estructura recurrente};
  \end{tikzpicture}
  \caption{Arquitectura del sistema protodeductivo. Cada regla transforma un
  patrón de premisas recurrentes en una afirmación de estabilidad, una relación
  generativa compuesta o una equivalencia estructural.}
  \label{fig:sistema-protodeductivo}
\end{figure}

\FloatBarrier

\section{Conclusión y alcance formal}

La teoría sitúa la autorrecurrencia en el nivel fundacional y presenta la
estabilidad, la medida, la cercanía y la autosimilitud como fenómenos derivados.
La notación anterior separa la igualdad literal $=$ de la equivalencia recurrente
$\peq$, y distingue una forma de sus iteraciones. Con ello se evitan varias
ambigüedades de la formulación inicial.

Para convertir la propuesta en una teoría matemática completa todavía se debe
especificar un modelo de $\F$, definir rigurosamente $\peq$ y $\pentails$, y
demostrar la consistencia e independencia de los axiomas. Hasta entonces, el
sistema debe entenderse como un marco protoaxiomático: una organización formal
de intuiciones sobre la emergencia de estructura mediante autorrecurrencia.

% -----------------------------------------------------------------------------
% Capítulo I
% -----------------------------------------------------------------------------

\capitulo{Capítulo I}{Fundamentos de la matemática autorrecursiva}

\renewcommand{\thesection}{I.\arabic{section}}
\renewcommand{\theHsection}{capI.\arabic{section}}
\renewcommand{\theHequation}{capI.\arabic{section}.\arabic{equation}}
\setcounter{section}{0}
\setcounter{axiomacap}{-1}
\setcounter{defcap}{0}

\section{Protoorigen: la unidad raíz}

Toda teoría debe precisar qué admite como existente. En el presente sistema no
se introduce ninguna entidad interna que no pueda obtenerse por necesidad
estructural a partir del fundamento autorrecursivo.

\begin{axiomacap}[Raíz]
Existe una única entidad original $\raiz$, indivisible y sin exterior dentro de
la teoría, cuya existencia consiste en su propia regeneración:
\begin{equation}
  \raiz(\raiz)\peq\raiz.
  \label{eq:cap-raiz}
\end{equation}
La unicidad se refiere a su función como origen generativo.
\end{axiomacap}

La expresión \eqref{eq:cap-raiz} no identifica dos cadenas tipográficas, sino
que afirma que la autoaplicación de la raíz reproduce su tipo estructural. Por
ello se utiliza la equivalencia recurrente $\peq$ en lugar de una igualdad
literal.

\section{Naturaleza del fundamento}

\begin{defcap}[Autoforma]
Una \emph{autoforma} es una forma $\alpha\in\F$ para la cual está definida la
secuencia completa de autoaplicaciones
\[
  \alpha,
  \quad \alpha(\alpha),
  \quad \alpha\bigl(\alpha(\alpha)\bigr),
  \quad\ldots
\]
o, de manera abreviada, $\bigl(\iter{\alpha}{n}\bigr)_{n\in\N}$.
\end{defcap}

Una autoforma no se presupone como número, conjunto o valor. Es, en primera
instancia, una forma capaz de actuar recurrentemente sobre su propio resultado.
La estabilidad de esa acción se estudiará después; por tanto, ser autoforma no
equivale todavía a ser un punto fijo.

\begin{axiomacap}[Pureza formal]
El dominio interno de la teoría contiene solamente autoformas y las relaciones
generadas necesariamente por sus recurrencias. El espacio, el tiempo, la
cantidad, la medida y la geometría no se toman como datos primitivos.
\end{axiomacap}

Este axioma expresa una restricción del lenguaje interno. La metateoría sigue
utilizando lógica, conjuntos e índices naturales para formular y analizar el
sistema.

\section{Primeros estados emergentes}

La primera autoaplicación de $\alpha$ puede producir distintos comportamientos.
Con la notación ya establecida, se distinguen:
\begin{align}
  \text{persistencia:}
  &\quad \iter{\alpha}{1}\peq\iter{\alpha}{0},
  \label{eq:persistencia}
  \\
  \text{expansión:}
  &\quad \iter{\alpha}{1}\not\peq\iter{\alpha}{0},
  \label{eq:expansion}
  \\
  \text{degeneración:}
  &\quad \iter{\alpha}{1}=\inestable.
  \label{eq:degeneracion}
\end{align}
El símbolo $\inestable$ designa una forma cuya recurrencia no puede continuar o
no produce una estructura observable. No se identifica con el conjunto vacío
$\varnothing$.

\begin{defcap}[Estado protoestructural]
Una autoforma $\alpha$ alcanza un estado protoestructural cuando dos etapas
consecutivas de su órbita son recurrentemente equivalentes. En el primer nivel
no trivial, esta condición se escribe
\begin{equation}
  \iter{\alpha}{2}\peq\iter{\alpha}{1}.
  \label{eq:estado-protoestructural}
\end{equation}
\end{defcap}

La condición anterior representa autosimilitud persistente, pero no exige que
los dos estados sean literalmente idénticos.

\begin{axiomacap}[Suficiencia generativa]
Toda autoforma cuya órbita alcanza un patrón estable genera una protoestructura
válida. Las nociones posteriores de dimensión, transformación y medida deberán
definirse a partir de dicho patrón.
\end{axiomacap}

\section{Emergencia de la protosecuencia}

La autoaplicación reiterada de $\alpha$ produce la \emph{protosecuencia}
\begin{equation}
  \iter{\alpha}{0}=\alpha,
  \qquad
  \iter{\alpha}{n+1}
  =\alpha\bigl(\iter{\alpha}{n}\bigr),
  \qquad n\in\N,
  \label{eq:protosecuencia-cap}
\end{equation}
que coincide con la convención de \eqref{eq:iteracion}. Cada término registra
un nivel de transformación; la secuencia se \emph{pliega} cuando un estado
estructural anterior reaparece.

\begin{defcap}[Recurrencia plegada]
Una autoforma $\alpha$ es \emph{plegable} si existen enteros $n>m\geq0$ tales
que $\iter{\alpha}{n}\peq\iter{\alpha}{m}$. Su grado de pliegue es el primer
índice en el que ocurre una repetición:
\begin{equation}
  \gr(\alpha)
  =\min\Bigl\{n\in\N_{>0}:\exists m<n,
  \ \iter{\alpha}{n}\peq\iter{\alpha}{m}\Bigr\}.
  \label{eq:grado-pliegue}
\end{equation}
Si no hay repetición, se define $\gr(\alpha)=\infty$.
\end{defcap}

El grado $\gr(\alpha)$ localiza la primera repetición, mientras que
$\mathcal D(\alpha)$, definida en \eqref{eq:dimension}, mide el período estable
eventual. Estas cantidades pueden coincidir, pero describen aspectos distintos.

\begin{axiomacap}[Pliegue finito]
Las estructuras discernibles desde el interior de la teoría proceden de
autoformas plegables. Una órbita infinitamente divergente puede existir en la
metateoría, pero no constituye una estructura internamente observable.
\end{axiomacap}

\section{Generación de protoestructuras}

\begin{defcap}[Protoestructura generada]
Si $\alpha$ es plegable, su protoestructura es la órbita ordenada hasta la
primera repetición:
\begin{equation}
  \Proto{\alpha}
  =\bigl(
    \iter{\alpha}{0},
    \iter{\alpha}{1},
    \ldots,
    \iter{\alpha}{\gr(\alpha)-1}
  \bigr),
  \label{eq:protoestructura-generada}
\end{equation}
junto con la identificación recurrente que cierra el ciclo.
\end{defcap}

Se emplea una sucesión ordenada, y no un conjunto ordinario, porque el orden de
las iteraciones contiene información dinámica. La protoestructura de la raíz se
denota por $\Proto{\raiz}$.

La órbita estricta de una raíz persistente puede contener una sola clase de
equivalencia. Para expresar que de ella emerge una diversidad de formas se
introduce, por tanto, su cierre generativo:
\begin{equation}
  \Cl_r(\raiz)
  =\bigl\{\alpha\in\F:\raiz\vdash_r^{*}\alpha\bigr\},
  \label{eq:cierre-generativo}
\end{equation}
donde $\vdash_r^{*}$ es la clausura reflexiva y transitiva de la relación de
generación recurrente.

\begin{axiomacap}[Encierro emergente]
Toda estructura interna de la teoría debe generarse a partir de la raíz; en
particular, sus formas pertenecen a $\Cl_r(\raiz)$. No se admiten como primitivas
estructuras externas o predeterminadas.
\end{axiomacap}

Esta formulación sustituye la expresión informal ``subconjunto de las iteraciones
de la raíz''. Si se interpretara esa frase literalmente, una raíz que satisface
\eqref{eq:cap-raiz} no podría producir diversidad estructural.

\section{Protoidentidad y distinguibilidad}

\begin{defcap}[Equivalencia recursiva]
Dos autoformas plegables $\alpha$ y $\beta$ son recurrentemente equivalentes si
sus protoestructuras tienen el mismo patrón de transición y de cierre:
\begin{equation}
  \alpha\peq\beta
  \quad\Longleftrightarrow\quad
  \Proto{\alpha}\cong_r\Proto{\beta},
  \label{eq:igualdad-recursiva}
\end{equation}
donde $\cong_r$ denota un isomorfismo que preserva el orden recurrente.
\end{defcap}

La igualdad literal $=$ queda reservada para la identidad de objetos dentro de
la metateoría. La relación $\peq$ compara comportamientos recurrentes, de modo
que dos formas distintas pueden representar la misma estructura emergente.

\begin{axiomacap}[Equivalencia emergente]
La comparación interna entre autoformas depende exclusivamente de sus patrones
de recurrencia. En consecuencia, las propiedades observables deben ser
invariantes bajo $\peq$.
\end{axiomacap}

\section{Protoestabilidad y comportamiento dinámico}

Dentro del dominio observable, una autoforma plegable puede clasificarse por el
destino de su recurrencia:
\begin{description}[leftmargin=2.2em,style=nextline]
  \item[Estática.]
  $\iter{\alpha}{1}\peq\iter{\alpha}{0}$; la forma alcanza inmediatamente un
  punto fijo estructural.

  \item[Cíclica.]
  Existen $n>m\geq0$ tales que
  $\iter{\alpha}{n}\peq\iter{\alpha}{m}$, con período estable mayor que uno.

  \item[Degenerada.]
  Existe $n\in\N$ tal que $\iter{\alpha}{n}=\inestable$; la recurrencia deja de
  producir estados estructuralmente observables.
\end{description}

Una autoforma no plegable constituye además un comportamiento
\emph{divergente}; por el axioma de pliegue finito, queda fuera de la taxonomía
interna de estructuras discernibles. Los términos ``vive'', ``muere'' y
``genera'' pueden emplearse como metáforas, pero las tres condiciones anteriores
son sus definiciones matemáticas.

\section{Resumen de la notación del capítulo}

\begin{center}
\renewcommand{\arraystretch}{1.25}
\begin{tabular}{>{$}c<{$} p{0.66\textwidth}}
\toprule
\text{Símbolo} & \textbf{Interpretación} \\
\midrule
\raiz & Raíz autorrecursiva distinguida. \\
\alpha,\beta,\gamma & Autoformas genéricas. \\
\iter{\alpha}{n} & Iteración autorrecursiva de nivel $n$. \\
\Proto{\alpha} & Protoestructura ordenada generada por $\alpha$. \\
\gr(\alpha) & Primer índice de pliegue de $\alpha$. \\
\alpha\peq\beta & Equivalencia de patrones recurrentes. \\
\inestable & Estado inestable o recurrencia interrumpida. \\
\Cl_r(\raiz) & Cierre de formas generadas desde la raíz. \\
\bottomrule
\end{tabular}
\end{center}

\section{Conclusión del capítulo}

Los seis axiomas numerados del $0$ al $5$ ---el axioma raíz y cinco axiomas
adicionales--- establecen un primer núcleo fundacional. A partir de ellos se
obtienen una secuencia de autoaplicaciones, un criterio de pliegue, una
protoestructura ordenada, una equivalencia recurrente y una clasificación de
los comportamientos observables.

El capítulo no demuestra todavía que exista un modelo no trivial que satisfaga
simultáneamente todos los axiomas. Su aportación es precisar las condiciones que
deberá cumplir dicho modelo para que operaciones, geometrías, medidas y
relaciones puedan emerger sin ser introducidas como primitivas externas.

\end{document}
